% ============================================================================
%  Part II: TesseraID Specification (normative)
% ============================================================================
\section{TesseraID Specification}
\label{sec:spec}

This part is the normative reference. Implementations that conform to this specification produce and consume identifiers with the byte layout, semantics, and encoding rules described below. Where this document uses RFC 2119 keywords (MUST, SHOULD, MAY), they carry the standard meanings.

\subsection{Wire format}

A TesseraID is exactly 41 bytes (328 bits), structured as four contiguous fields:

\begin{center}
\begin{tikzpicture}[
    field/.style={draw,minimum height=14mm,inner sep=4pt,align=center,font=\small},
    label/.style={font=\scriptsize\itshape,gray},
    bytelabel/.style={font=\scriptsize,below=1mm}
  ]
  % Field rectangles (widths roughly proportional to bit widths)
  \node[field, fill=blue!10,    minimum width=24mm] (ts)  at (0,0)    {\textbf{Timestamp}\\48 bits\\6 bytes};
  \node[field, fill=orange!15,  minimum width=10mm, right=0pt of ts]  (ver) {\textbf{Variant}\\8 bits\\1 byte};
  \node[field, fill=green!15,   minimum width=14mm, right=0pt of ver] (ctr) {\textbf{Counter}\\16 bits\\2 bytes};
  \node[field, fill=purple!10,  minimum width=84mm, right=0pt of ctr] (rnd) {\textbf{Random tail}\\256 bits\\32 bytes};

  % Byte offset markers
  \node[bytelabel] at (ts.south west) {0};
  \node[bytelabel] at (ts.south east) {6};
  \node[bytelabel] at (ver.south east) {7};
  \node[bytelabel] at (ctr.south east) {9};
  \node[bytelabel] at (rnd.south east) {41};

  % Total span brace
  \draw[decorate,decoration={brace,amplitude=4pt}, gray]
    ($(ts.north west)+(0,2mm)$) -- ($(rnd.north east)+(0,2mm)$)
    node[midway,above=4pt,font=\small\itshape] {41 bytes total --- 55 chars in Base64URL --- 66 chars in Crockford Base32};
\end{tikzpicture}
\end{center}

\textbf{Field semantics}:
\begin{description}
  \item[Timestamp (offsets 0--5).] Big-endian unsigned 48-bit integer, milliseconds since the Unix epoch (1970-01-01T00:00:00Z). The 48-bit field provides $2^{48}$ ms $\approx$ 8{,}919 years of forward range, sufficient through the year 10889.
  \item[Variant (offset 6).] \label{sec:variant-reservation} Single byte.
    \begin{itemize}
      \item \code{0x01} = Standard (wall-clock).
      \item \code{0x02} = HLC (Hybrid Logical Clock).
      \item \code{0x03} = TesseraHash (sibling content-identifier format, §\ref{sec:tessera_hash}). Allocated 2026 from the reserved-future range; TesseraID decoders MUST reject \code{0x03} (a 0x03 buffer is not a TesseraID).
      \item \code{0x00} and \code{0xFF} are \textbf{permanently reserved}: they MUST be rejected by conformant decoders and MUST NOT be assigned in any future revision of this specification. These two values are reserved defensively so that TesseraID buffers cannot be confused with zero-filled or all-ones uninitialized memory.
      \item \code{0x04} = \tesseraseal{} (sibling encryption-envelope format; allocated 2026 from the reserved-future range). TesseraID decoders MUST reject \code{0x04} with an unknown-variant error: a \code{0x04} buffer is a \tesseraseal{} envelope, not a TesseraID.
      \item \code{0x05}--\code{0xFE} are reserved for future variants. Conformant decoders MUST reject them with an unknown-variant error until they are assigned by a future revision of this specification.
    \end{itemize}
  \item[Counter (offsets 7--8).] Big-endian unsigned 16-bit integer. For Standard, the counter is incremented within a single millisecond by a small random step (1--16); for HLC, it is the HLC logical counter.
  \item[Random tail (offsets 9--40).] 32 bytes from a cryptographically-strong random source. The full 256 bits provide $2^{128}$ effective post-quantum security against second-preimage attacks.
\end{description}

\subsection{Standard variant generation (variant \code{0x01})}

The Standard generation algorithm assumes a wall-clock time source that returns milliseconds since the Unix epoch. It maintains two pieces of per-generator state: \code{last\_ts} (initially 0) and \code{counter} (initially 0).

\begin{lstlisting}[language=Pseudocode,caption={Standard TesseraID generation},label={lst:gen-std}]
procedure new_standard():
    rnd[0..32]  <-  csprng()           # 33 bytes total
    ts          <-  now_ms()
    if ts < last_ts:                   # clamp on regression
        ts  <-  last_ts
    if ts == last_ts:                  # same ms as previous ID
        step    <-  1 + (rnd[0] & 0x0F)    # random step in [1, 16]
        next    <-  counter + step
        if next > 0xFFFF:              # counter exhausted, spin
            while now_ms() <= last_ts: pass
            ts       <-  now_ms()
            counter  <-  rnd[0]        # reseed at small value
        else:
            counter  <-  next
    else:                              # new ms
        counter  <-  rnd[0]            # randomized initial offset
    last_ts <-  ts
    return  be48(ts) || 0x01 || be16(counter) || rnd[1..32]
\end{lstlisting}

The random step (rather than a fixed +1 increment) is a deliberate choice. It provides modest unlinkability of consecutive IDs from a single generator while preserving strict monotonicity (the step is always at least 1, so consecutive IDs always increase).

\emph{Non-normative remark on the exhaustion spin:} the spin condition \code{while now\_ms() <= last\_ts} is deliberately \code{<=}, not \code{==}. After a clock regression, \code{last\_ts} may sit ahead of the wall clock; spinning until \code{now\_ms()} strictly exceeds \code{last\_ts} guarantees the post-spin timestamp is strictly greater than every previously emitted timestamp, preserving strict monotonicity even when the regression spans the exhaustion event. An implementation that spins only while \code{now\_ms() == last\_ts} exits the loop immediately during a regression and can emit an identifier that sorts before its predecessor.

\subsection{Optional defensive helper: \code{validate\_min\_timestamp}}
\label{sec:min-timestamp}

Implementations SHOULD expose a non-normative helper that returns true iff a decoded TesseraID's timestamp is at or after a project epoch \code{MIN\_TIMESTAMP\_MS}. The reference value is \code{1{,}700{,}000{,}000{,}000} ms ($\approx$ 2023-11-14), chosen as the earliest plausible timestamp any TesseraID-aware deployment will ever emit.

This helper is \textbf{not part of the wire contract}. Standard decoders MUST continue to accept all timestamps in $[0, 2^{48} - 1]$ as required by the field semantics in §\ref{sec:why}. The helper exists for callers who want to cheaply reject obvious mis-parses — a 41-byte buffer that decodes to timestamp~0 has either come from zeroed memory, a wrong-offset read, or a buggy producer. Rejecting such inputs at boundary entry points is a defense-in-depth measure, not a conformance requirement.

\subsection{CSPRNG requirements}
\label{sec:csprng}

The 256-bit random tail is the only field carrying the post-quantum security claim. Conformant implementations MUST therefore use an operating-system-provided cryptographically-strong pseudorandom number generator. The following normative rules apply:

\begin{description}
  \item[Source.] Generators MUST draw the random tail from one of: \code{getrandom(2)} on Linux 3.17+, \code{arc4random\_buf} or \code{getentropy(3)} on the BSDs and macOS, \code{BCryptGenRandom} on Windows, the language-level wrappers around these (\code{os.urandom}, \code{crypto/rand.Read}, \code{getrandom::fill}, \code{crypto.getRandomValues}, \code{SecureRandom}), or — as a last-resort portable fallback — \code{/dev/urandom}. Generators MUST NOT use a non-cryptographic userspace PRNG (Mersenne Twister, Xoshiro, \code{rand}, etc.) under any circumstance.
  \item[Failure.] If the CSPRNG returns an error or refuses to produce 32 bytes, the generator MUST surface the error to the caller and MUST NOT emit a TesseraID. There is no fallback path.
  \item[Partial reads.] Partial reads (an underlying \code{read()} or \code{getrandom()} returning fewer than 32 bytes) MUST be retried until either 32 bytes are obtained or an unrecoverable error is returned.
  \item[Independence.] Each TesseraID's random tail MUST be drawn independently. A generator MAY combine the 32-byte tail with a small additional read (1--2 bytes) used for the counter step in the same syscall as an optimization; this is unrelated to the cryptographic claim and does not weaken it because the OS CSPRNG output is uniform.
\end{description}

\noindent A test vector cannot directly verify entropy quality, but conformance reviewers SHOULD spot-check each implementation's call site against this list. The current implementations call: \code{os.urandom} (Python), \code{getrandom}/\code{getentropy} (C/C++), \code{crypto/rand.Read} (Go), \code{getrandom::fill} (Rust), \code{crypto.getRandomValues}/\code{randomFillSync} (TypeScript), \code{SecureRandom} (Java), and the underlying C reference (CPython extension).

\subsection{HLC variant generation (variant \code{0x02})}

The HLC generator maintains \code{l} (logical time, ms-grained) and \code{c} (logical counter). On every generation, it consults both the local physical clock and any peer HLC observation.

The contract uses two thresholds, distinct in purpose:

\begin{description}
  \item[Hard outer bound $\Delta_{\max} = 60{,}000$ ms (normative).] Conformant generators MUST reject observations whose logical time exceeds the local physical clock by more than $\Delta_{\max}$. The comparison is \emph{strict}: an observation MUST be rejected iff $l_{\mathit{msg}} - pt > 60{,}000$ ms --- an observation exactly $60{,}000$ ms ahead MUST be accepted, and one $60{,}001$ ms ahead MUST be rejected. A rejected observation MUST leave the generator's $(l, c)$ state unchanged. These boundary semantics are pinned by the \texttt{hlc\_drift} conformance vectors (TESS-00009). Implementations MUST expose this constant under the name \code{HLC\_MAX\_DRIFT\_MS} (or the language's idiomatic equivalent). $\Delta_{\max}$ is part of the wire contract: a 60-second peer-ahead-of-local observation that one implementation accepts and another rejects is an interoperability bug.
  \item[Local tightening override (optional).] $\Delta_{\max} = 60{,}000$ ms is the normative \emph{maximum} and the default. Implementations MAY expose a per-generator override that only \emph{tightens} the bound: a configured override MUST satisfy $1~\text{ms} \le \Delta < 60{,}000$~ms. Configuring a bound larger than $\Delta_{\max}$ MUST surface an error and produce no output (relaxing the bound defeats the monotonicity guarantee). Where the override API uses an integer field, the value $0$ selects the normative default --- matching the C reference's \code{max\_drift\_ms} semantics (TID-00020). Implementations that expose no override are conformant.
  \item[Soft alarm threshold $\Delta_{\text{alarm}}$ (configurable, default 1{,}000 ms).] Implementations SHOULD emit a log, metric, or alarm when an observation exceeds $\Delta_{\text{alarm}}$, even if it is still below $\Delta_{\max}$. The default of 1 second captures the operational reality that well-disciplined hosts (NTP, chrony, AWS Time Sync, GCP) maintain sub-100~ms drift; a peer 1~s ahead almost certainly indicates a clock-discipline or NTP-outage problem that an operator should know about, even though the wire format still accepts it. $\Delta_{\text{alarm}}$ is NOT part of the wire contract.
\end{description}

\noindent The rationale for $\Delta_{\max} = 60$~s: large enough to absorb a brief NTP outage, container time injection, or VM live-migration without breaking valid traffic; small enough that a maliciously advanced peer cannot stamp identifiers a year into the future. The HLC literature (Kulkarni et al.~2014, \S3) explicitly leaves the bound to operators; the CockroachDB and kronos operational practice in production is approximately 500~ms to 1~s. We pick 60~s for the hard bound because the wire-format invariant should tolerate the worst plausible operational event without changing, and we pick 1~s as the default alarm because that's where well-run systems already live.

\begin{lstlisting}[language=Pseudocode,caption={HLC TesseraID generation (with optional peer observation)},label={lst:gen-hlc}]
procedure new_hlc(observe = none):
    rnd[0..31] <- csprng()             # 32 bytes
    pt         <- now_ms()             # physical time
    l_old, c_old <- l, c
    if observe is not none:
        l_msg, c_msg <- observe
        if (l_msg > pt) and ((l_msg - pt) > MAX_DRIFT_MS):
            error  HLC_drift_exceeded
        l_new <- max(l_old, l_msg, pt)
        case:
            l_new == l_old == l_msg:  c_new <- max(c_old, c_msg) + 1
            l_new == l_old:           c_new <- c_old + 1
            l_new == l_msg:           c_new <- c_msg + 1
            default:                  c_new <- 0
    else:
        l_new <- max(l_old, pt)
        c_new <- (c_old + 1) when l_new == l_old, else 0
    if c_new > 0xFFFF: l_new += 1; c_new <- 0
    l, c <- l_new, c_new
    return  be48(l_new) || 0x02 || be16(c_new) || rnd
\end{lstlisting}

\subsection{Text encodings}

TesseraID defines two text encodings. Both round-trip cleanly to the same 41-byte binary form and neither uses padding. They differ in their behavior under string sort: Crockford Base32 \citep{crockford} is ASCII-sort-preserving (sorting the strings as text gives generation-time order), while canonical Base64URL is sort-preserving only when its bytes are compared, not when its text is compared as ASCII. The reason is that the RFC 4648 §5 alphabet \citep{rfc4648} (\code{A--Z}, \code{a--z}, \code{0--9}, \code{-}, \code{\_}) maps alphabet values 0..63 to ASCII codepoints that are not monotonic --- \code{A--Z} (65..90), \code{a--z} (97..122), \code{0--9} (48..57), \code{-} (45), \code{\_} (95) interleave, so a string-level ASCII sort produces an order different from the underlying byte order. Crockford's alphabet (\code{0--9} then \code{A--Z} minus \code{I/L/O/U}) is monotonic in ASCII and so its strings sort identically to their decoded bytes. Choose accordingly: machine-to-machine paths use Base64URL and sort the bytes (or simply preserve insertion order); human-facing or text-sorted columns use Crockford.

\subsubsection{Canonical: Base64URL (55 characters)}

RFC 4648 Section 5 (``Base 64 Encoding with URL and Filename Safe Alphabet''), no padding. Encoding 41 bytes produces $\lceil 41 \cdot 8 / 6 \rceil = 55$ characters; the final character carries 4 bits of value plus 2 bits of zero padding. The alphabet is \code{A-Z}, \code{a-z}, \code{0-9}, \code{-}, \code{\_}. Decoders MUST reject inputs whose final character carries non-zero padding bits.

\subsubsection{Human-typable: Crockford Base32 (66 characters)}

Crockford's alphabet (\code{0123456789ABCDEFGHJKMNPQRSTVWXYZ}) and its case-insensitive decoding (with the I/L→1 and O→0 confusables canonicalized) make this encoding suitable for human transcription. Encoding 41 bytes produces $\lceil 41 \cdot 8 / 5 \rceil = 66$ characters; the final character carries 3 bits of value plus 2 bits of zero padding.

The Crockford form is intended for use cases where humans handle the identifier directly: pasted into chat messages, read aloud over phones, written on paper, or appearing in CLI tools where mixed-case is hostile. The Base64URL form is intended for everything else.

\subsubsection{Decode-time variant validation (both encodings)}
\label{sec:decode-variant-contract}

Both text decoders --- the 55-character Base64URL decoder and the 66-character Crockford Base32 decoder --- MUST validate the variant byte of the decoded 41-byte value and MUST reject any value other than \code{0x01} or \code{0x02} with the unknown-variant error \emph{at decode time}. Variant validation is part of decode, not a separate post-decode step: a conformant text decoder never returns a 41-byte buffer whose variant byte it does not own.

\emph{Non-normative remark:} combined with the padding-bit guards above, decode-time variant validation makes each text decoder a bijection onto the set of valid TesseraIDs --- every accepted string maps to exactly one valid TesseraID, and every valid TesseraID maps to exactly one accepted canonical string.

\subsection{TesseraHash --- sibling content-identifier format (variant \code{0x03})}
\label{sec:tessera_hash}

TesseraHash is a \textbf{sibling format} to TesseraID, sharing TesseraID's 9-byte header layout but extending the wire size and replacing the random tail with a deterministic content digest. TesseraHash is \emph{additive}: it does not change TesseraID, and the two formats are guaranteed to not cross-decode.

\emph{This subsection is a non-normative summary for TesseraID implementers.} The normative source for the TesseraHash wire format, its encodings, and its security claims is the TesseraHash whitepaper (LaTeX sources under \texttt{tessera-hash/latex/}); where this summary and that paper disagree, the TesseraHash paper wins. The only normative statement TesseraID makes about TesseraHash is the cross-format rejection contract in §\ref{sec:variant-reservation} and below.

\subsubsection{Wire layout (73 bytes)}

\begin{lstlisting}
0       6      7       9                                            73
+-------+------+-------+--------------------------------------------+
|  TS   | 0x03 |  CTR  |              SHAKE256(content, 64B)        |
| 48 b  |  8 b | 16 b  |                       512 b                |
+-------+------+-------+--------------------------------------------+
\end{lstlisting}

\begin{description}
  \item[Timestamp (offsets 0--5).] Big-endian unsigned 48-bit ms since Unix epoch. Provenance metadata: \emph{when} this identifier was assigned. NOT a security claim.
  \item[Variant (offset 6).] Fixed at \code{0x03}, allocated from TesseraID's reserved-future range (§\ref{sec:variant-reservation}).
  \item[Counter (offsets 7--8).] Big-endian unsigned 16-bit. Same-ms disambiguator within a generator.
  \item[Tail (offsets 9--72).] 64 bytes (512 bits) of \code{SHAKE256(content)} \citep{fips202}, output truncated to 64 bytes. The tail is \emph{deterministic from the content alone}: two callers that TesseraHash the same content produce the same tail, regardless of generator or timestamp.
\end{description}

\subsubsection{Security claim}

The 512-bit tail gives:
\begin{itemize}
  \item Second-preimage resistance of $\min(2^d, 2^{256})$ per FIPS 202 \citep{fips202} --- $2^{256}$ classical at $d = 512$ output bits (the SHAKE-256 capacity cap, not the output width, is the binding term). The generic quantum (Grover \citep{grover1996}) bound on the same search is also $2^{256}$, so the figure holds post-quantum.
  \item $2^{256}$ classical collision resistance (birthday).
  \item $2^{170}$ post-quantum collision resistance (BHT \citep{bht1998}). This is the surface that matters for content addressing across multi-decade archival horizons.
\end{itemize}

\textbf{What TesseraHash is NOT:} not a hash function (use SHAKE-256 directly), not a MAC (use KMAC256), not a signature (use ML-DSA / SLH-DSA). TesseraHash is a content identifier with provenance metadata; the timestamp carries no security claim. See §\ref{sec:threat-model} for the framing TesseraID established; the same boundary applies, sharpened: the second-preimage figure for TesseraHash is the brute-force-the-content-via-SHAKE game, which is what every content-addressed format gets.

\subsubsection{Cross-format non-overlap}

TesseraID and TesseraHash MUST NOT cross-decode:
\begin{itemize}
  \item TesseraID is exactly 41 bytes with variant \code{0x01} or \code{0x02}. Conformant TesseraID decoders reject 73-byte buffers (length error) and reject the \code{0x03} variant byte (unknown-variant error, codified by the §\ref{sec:variant-reservation} reservation).
  \item TesseraHash is exactly 73 bytes with variant \code{0x03}. Conformant TesseraHash decoders reject 41-byte buffers (length error) and reject any variant other than \code{0x03}.
\end{itemize}

An application that may receive either format dispatches at the application boundary:

\begin{lstlisting}
if len(buf) == 73 and buf[6] == 0x03:  TesseraHash.decode(buf)
elif len(buf) == 41 and buf[6] in {0x01, 0x02}:  TesseraID.decode(buf)
else: error
\end{lstlisting}

\subsubsection{Encodings}

TesseraHash uses the same Base64URL and Crockford Base32 alphabets as TesseraID, sized for the larger payload:

\begin{itemize}
  \item \textbf{Base64URL: 98 characters.} $\lceil 73 \cdot 8 / 6 \rceil = 98$. 584 input bits fill 97 full characters plus 2 bits, so the final character carries 2 bits of value plus 4 bits of zero padding (it is always one of \code{A}, \code{Q}, \code{g}, \code{w}). Decoders MUST reject inputs whose final character carries non-zero padding bits.
  \item \textbf{Crockford Base32: 117 characters.} $\lceil 73 \cdot 8 / 5 \rceil = 117$. 584 bits of input produce 585 output bits with 1 zero pad bit at the bottom of the final character. Case-insensitive decoding with the same I/L→1 and O→0 confusables as TesseraID.
\end{itemize}

\subsubsection{Verify}

A conformant TesseraHash implementation MUST expose a \code{verify(raw, content)} predicate that returns true iff \code{raw} is a 73-byte TesseraHash with variant \code{0x03} \emph{and} its tail equals \code{SHAKE256(content)} truncated to 64 bytes. The tail comparison MUST be constant-time on the 64-byte input to avoid timing oracles. Malformed buffers return false rather than raising.

\subsubsection{Generator semantics}

The TesseraHash generator maintains \code{last\_ts} and \code{counter}. On every call:
\begin{enumerate}
  \item Compute \code{tail = SHAKE256(content, 64)}.
  \item Read the wall clock; clamp on regression (as for TesseraID Standard).
  \item If the new timestamp equals \code{last\_ts}, increment the counter (spin to next ms on overflow).
  \item Otherwise reset the counter to 0.
\end{enumerate}

Consecutive calls from a single generator produce strictly monotonic identifiers (binary order, and ASCII order under Crockford). Two callers that TesseraHash the same content at different times produce identifiers that differ only in \code{(ts, counter)} and share the 64-byte tail.

